\documentclass[12pt,a4paper]{article}
\usepackage[utf8]{inputenc}
\usepackage{amsmath}
\usepackage{amsfonts}
\usepackage{amssymb}
\usepackage{amsthm}
\usepackage{hyperref}

\usepackage{multicol}
\usepackage{fancyhdr}
\usepackage[inline]{enumitem}
\usepackage{tikz}
\usepackage{tikz-cd}
\usetikzlibrary{calc}
\usetikzlibrary{shapes.geometric}
\usepackage[margin=0.5in]{geometry}
\usepackage{xcolor}
\usepackage{listings} 

\hypersetup{
    colorlinks=true,
    linkcolor=blue,
    filecolor=magenta,      
    urlcolor=cyan,
    pdftitle={Tensors},
    pdfpagemode=FullScreen,
    }

%\pagestyle{fancy}
% 	\fancyhf{}
%\chead{ \fancyplain{}{\CLASSNAME} }
%\chead{ \fancyplain{}{\STUDENTNAME} }
\rhead{\thepage}
\input{../MyMacros}

\newcommand{\RED}[1]{\textcolor{red}{*1}}
\newcommand{\BLUE}[1]{\textcolor{blue}{*1}}
\definecolor{mygreen}{RGB}{28,172,0} % color values Red, Green, Blue
\definecolor{mylilas}{RGB}{170,55,241}

\begin{document}

\begin{center}
	\textbf{\Large{High Order Weighted
Essentially Nonoscillatory
Schemes for Convection
Dominated Problems}}
\end{center}
\begin{center}
	A review by Paul Carmody\\Math 5315 Numerical Methods for PDE: Fall 2025
\end{center}

\HLINE
\begin{abstract}
	Provide a synapses and understanding of the article ``High Order Weighted
Essentially Nonoscillatory
Schemes for Convection
Dominated Problems" by Chi-Wang Shu.  The background and motivation for the WENO process is described along with its foundation in ENO (and its limitations).  The WENO methodology is then described in a step by step manner.  We measure than the behavior of the non-linear weights (decribed below).  This article ends with a discussion measuring the merits, advantages and disadvantages of this approach.
\end{abstract}

\HLINE
The Weighted Essentially Non-Oscillatory (WENO) scheme is a high-resolution numerical method widely used for solving hyperbolic partial differential equations, particularly those involving discontinuities, shocks, or sharp gradients. It is designed to achieve high-order accuracy in smooth regions while simultaneously suppressing spurious oscillations near discontinuities, a common problem with standard high-order schemes (like central difference schemes).\\

\noindent Reasoning and Explanation

\begin{description}

\item Background and Motivation

Hyperbolic conservation laws often take the form:
$$\frac{\partial u}{\partial t} + \frac{\partial f(u)}{\partial x} = 0$$
where $u$ is the conserved quantity and $f(u)$ is the flux function.

When solving these equations numerically, especially using finite difference or finite volume methods, achieving high accuracy (order $p \ge 3$) is desirable. However, Godunov's theorem (or the TVD theorem) states that linear schemes that are Total Variation Diminishing (TVD) must be at most first-order accurate near extrema. High-order linear schemes (like central differences) introduce oscillations (Gibbs phenomena)\cite{gibbs} near discontinuities.

Non-linear schemes, such as TVD schemes (e.g., MUSCL), overcome this limitation but often cap the accuracy at second order. The WENO scheme is a non-linear approach that achieves high-order accuracy (typically fifth order) in smooth regions while maintaining non-oscillatory properties near shocks.

\item  The ENO Foundation

WENO is an improvement upon the Essentially Non-Oscillatory (ENO)\cite{eno} scheme, developed by Harten, Engquist, Osher, and Chakravarthy.

\begin{description}

\item ENO Principle: 

The ENO scheme achieves non-oscillatory behavior by adaptively choosing the smoothest stencil (a set of neighboring grid points) from a collection of candidate stencils to construct the numerical flux or interpolation polynomial. If a stencil spans a discontinuity, it is discarded in favor of a smoother one.

\item Limitation of ENO:

 ENO is only "essentially" non-oscillatory, meaning the error is $O(\Delta x^p)$ in smooth regions, but the scheme is not optimally accurate. The selection process is discrete (either choose stencil A or stencil B), which leads to a non-smooth flux function and a slight reduction in convergence rate compared to its potential order.

\end{description}
\item The WENO Methodology

The WENO scheme, introduced by Liu, Osher, and Chan \cite{weno-l} overcomes the limitations of ENO by using a \textit{convex combination} of all candidate stencils, weighted by non-linear weights.

\begin{description}

\item  Step 1: Flux Splitting (Required for Hyperbolic Systems)

To ensure stability and handle the direction of information propagation, the flux function $f(u)$ is typically split into positive and negative parts based on the sign of the eigenvalues of the Jacobian $\partial f / \partial u$. This is often done using the Lax-Friedrichs splitting or characteristic decomposition:
$$f(u) = f^+(u) + f^-(u)$$
The numerical flux $\hat{f}_{i+1/2}$ is then calculated as the sum of the fluxes derived from $f^+$ and $f^-$. We focus on calculating the numerical flux for $f^+$ (the procedure is symmetric for $f^-$).

\item Step 2: Candidate Stencils and Interpolation

For a $p$-th order WENO scheme (e.g., $p=5$), we define $r$ small stencils, $S_k$, each containing $r$ points, such that their union forms a large stencil $S$ of size $2r-1$. For 5th-order WENO ($r=3$), we use three 3-point stencils:

\begin{tabular}{|c|c|c|c|}
\hline
	Stencil Index ($k$) & Stencil Points & &\\
	\hline
	0 &  $u_{i-2}$& $u_{i-1}$& $u_i$ \\
	\hline 
	1 & 	$u_{i-1}$ & $ u_i$& $u_{i+1}$ \\
	\hline
	2 & $u_i$ &$u_{i+1}$ & $ u_{i+2}$ \\
	\hline 
\end{tabular}

For each small stencil $S_k$, a high-order polynomial interpolation $P_k(x)$ is constructed to approximate the flux function $f(u)$ at the interface $x_{i+1/2}$. This yields a candidate numerical flux $\hat{f}_{i+1/2}^{(k)}$.

\item Step 3: Smoothness Indicators ($\beta_k$)

The core of the WENO scheme is the determination of how smooth the solution is within each candidate stencil $S_k$. This is quantified by a smoothness indicator, $\beta_k$. $\beta_k$ is typically a measure of the $L_2$ norm of the derivatives of the interpolation polynomial $P_k(x)$ over the stencil.

In practice, $\beta_k$ is calculated using discrete differences of the cell values $u_i$. A smaller $\beta_k$ indicates a smoother solution within stencil $S_k$. If a discontinuity lies within $S_k$, $\beta_k$ will be large.

\item Step 5: Non-linear Weights ($\omega_k$)

The non-linear weights $\omega_k$ are calculated based on the linear weights $d_k$ and the smoothness indicators $\beta_k$. The goal is to assign a weight close to $d_k$ if the stencil is smooth, and a weight close to zero if the stencil contains a discontinuity.

Behavior of Non-linear Weights:

\begin{description}

\item  Smooth Region: 

If all stencils are smooth, $\beta_k \approx O(\Delta x^2)$. The weights $\omega_k$ approach the linear weights $d_k$, and the scheme achieves the full high order (e.g., 5th order).

\item  Discontinuity Present:

 If stencil $S_j$ contains a discontinuity, $\beta_j$ will be large, causing $\alpha_j$ and consequently $\omega_j$ to approach zero. The scheme effectively switches to a lower-order, non-oscillatory interpolation based on the remaining smooth stencils.

\end{description}

\item Step 6: Final Numerical Flux

The final numerical flux $\hat{f}_{i+1/2}$ is the weighted average of the candidate fluxes:
$$\hat{f}_{i+1/2} = \sum_{k=0}^{r-1} \omega_k \hat{f}_{i+1/2}^{(k)}$$

This flux is then used in the semi-discrete finite volume formulation:
$$\frac{d u_i}{d t} = -\frac{1}{\Delta x} (\hat{f}_{i+1/2} - \hat{f}_{i-1/2})$$

\end{description}

\item 4. Advantages and Disadvantages

\begin{description}

\item High Accuracy

Achieves high order (typically 5th or 7th) in smooth regions.  

\item Non-Oscillatory

Effectively suppresses spurious oscillations near shocks and discontinuities. 

\item Robustness

Highly robust for complex flow problems, including turbulence and compressible flows. 

\item Complexity

Significantly more complex to implement than linear schemes due to flux splitting, multiple stencil calculations, and non-linear weighting.  

\item Computational Cost

Higher computational cost per time step compared to lower-order schemes (e.g., 2nd order TVD) due to the multiple flux evaluations and weight calculations.

\item Accuracy Near Extrema

While high-order in smooth regions, the scheme typically drops to 3rd order accuracy near critical points (smooth extrema) due to the non-linear weighting function. 

\end{description}



\item In Summary

The Weighted Essentially Non-Oscillatory (WENO) scheme is a high-resolution numerical method for solving hyperbolic PDEs. It achieves high-order accuracy (e.g., 5th order) in smooth regions while maintaining non-oscillatory behavior near discontinuities (shocks).

The core mechanism involves:

\begin{enumerate}

\item  Flux Splitting:

Decomposing the flux into positive and negative components.

\item  Multiple Stencils:

Constructing candidate numerical fluxes from several small, overlapping stencils.

\item  Smoothness Indicators ($\beta_k$):

 Quantifying the smoothness of the solution within each stencil.

\end{enumerate}

\item Conclusion:

The Weighted Essentially Non-Oscillatory (WENO) scheme is a high-resolution method for numerically solving hyperbolic PDEs that achieves high-order accuracy in smooth regions and suppresses spurious oscillations near discontinuities by computing the numerical flux as an adaptive weighted average over several candidate stencils, with non-linear weights determined by smoothness indicators for each stencil.

\end{description}

\begin{thebibliography}{1}

\bibitem{eno}
A. Harten, B. Engquist, S. Osher, and S. Chakravarthy, Uniformly high order essentially non-oscillatory schemes, III, J. Comput. Phys., 71 (1987), pp. 231–303.

\bibitem{gibbs}
D. Gottlieb and C.-W. Shu, On the Gibbs phenomenon and its resolution, SIAM Rev., 39 (1997), pp. 644–668.

\bibitem{weno-l}
X.-D. Liu, S. Osher, and T. Chan, Weighted essentially non-oscillatory schemes, J. Comput. Phys., 115 (1994), pp. 200–212.
\end{thebibliography}
\end{document}
